\chapter{Gestational Diabetes}

\section*{Definition and Overview}
Gestational diabetes mellitus (GDM) is a type of diabetes that develops during pregnancy in women who did not previously have diabetes. It occurs when the placenta produces hormones that interfere with the action of insulin, so the body cannot keep blood glucose levels within the normal range. Gestational diabetes usually develops in the second half of pregnancy, is often detected by routine screening, and typically resolves after the baby is born. Managing gestational diabetes well, with diet, physical activity, and sometimes insulin or other medicines, keeps the mother and baby healthy and reduces the risk of complications during pregnancy and birth.

\section*{Epidemiology}
Gestational diabetes is common and becoming more so. In Australia, around 14--17\% of pregnancies are affected, and rates have risen sharply over recent decades in line with rising obesity and maternal age, and with more sensitive diagnostic criteria. GDM is more common in women who are overweight, older than 35, of certain ethnic backgrounds (including South Asian, Middle Eastern, and Pacific Islander), who have a family history of type 2 diabetes, or who have had gestational diabetes or a large baby in a previous pregnancy. Women with GDM have a substantially increased risk of developing type 2 diabetes later in life.

\section*{Aetiology and Pathophysiology}
Gestational diabetes arises from the interaction of pregnancy-related hormonal changes with an underlying susceptibility to insulin resistance.

\begin{itemize}
    \item \textbf{Placental hormones:} hormones such as human placental lactogen, cortisol, and progesterone rise through pregnancy and reduce the effectiveness of insulin
    \item \textbf{Insulin resistance:} the body's cells respond less effectively to insulin, so the pancreas must produce more insulin to control blood glucose
    \item \textbf{Beta-cell insufficiency:} in women who develop GDM, the pancreas cannot increase insulin production enough to overcome the resistance
    \item \textbf{Rising glucose:} as pregnancy progresses, blood glucose levels rise, usually becoming apparent from around 24--28 weeks
    \item \textbf{Placental transfer:} glucose crosses the placenta to the baby, stimulating the fetal pancreas to produce extra insulin
    \item \textbf{Fetal effects:} the extra glucose and insulin cause the baby to grow large (macrosomia) and can affect lung, metabolic, and blood-sugar regulation after birth
\end{itemize}

\section*{Clinical Features}
Gestational diabetes usually causes no symptoms and is detected by screening. Some women may notice:

\begin{itemize}
    \item Increased thirst and urination
    \item Fatigue and tiredness
    \item Blurred vision in some cases
    \item Recurrent thrush or urinary tract infections
    \item The baby measuring large for dates on ultrasound
    \item Most women have no symptoms at all, which is why routine screening is important
\end{itemize}

\section*{Differential Diagnosis}
When high blood glucose is found in pregnancy, the type of diabetes must be established.

\begin{itemize}
    \item Pre-existing type 1 or type 2 diabetes that was undiagnosed before pregnancy
    \item Gestational diabetes, diagnosed for the first time in pregnancy
    \item Transient glucose intolerance related to medicines such as corticosteroids
    \item Polycystic ovary syndrome with insulin resistance, which may coexist
    \item Other causes of glycosuria (glucose in the urine) without diabetes
\end{itemize}

\section*{When to Seek Medical Care}
All pregnant women in Australia are offered screening for gestational diabetes, usually with an oral glucose tolerance test (OGTT) at 24--28 weeks. Women with risk factors may be tested earlier. Anyone diagnosed with GDM needs regular antenatal, diabetes, and dietary follow-up.

\textbf{Call 000 (or local emergency services) for:}
\begin{itemize}
    \item Signs of very high blood glucose, including severe thirst, frequent urination, vomiting, or feeling confused and unwell
    \item Reduced fetal movements, which requires urgent obstetric review
    \item Symptoms of pre-eclampsia such as severe headache, visual disturbance, or upper abdominal pain
    \item Any emergency in pregnancy, including bleeding, severe pain, or waters breaking early
\end{itemize}

\section*{Diagnosis and Investigations}
Screening and diagnosis are based on blood glucose testing under standardised conditions.

\begin{itemize}
    \item \textbf{Screening:} all women are offered screening between 24 and 28 weeks of pregnancy
    \item \textbf{Oral glucose tolerance test (OGTT):} fasting blood glucose is measured, then a glucose drink is given and blood glucose is measured again at one and two hours
    \item \textbf{Diagnostic criteria:} GDM is diagnosed if any glucose level is at or above the agreed thresholds for the fasting, one-hour, or two-hour sample
    \item \textbf{Self-monitoring:} women diagnosed with GDM measure blood glucose levels at home with a finger-prick meter, typically before and after meals
    \item \textbf{Monitoring the baby:} ultrasound scans track fetal growth and amniotic fluid
    \item \textbf{Postnatal testing:} an OGTT is repeated at six weeks after birth to confirm that glucose levels have returned to normal
\end{itemize}

\section*{Management}
The goal of treatment is to keep blood glucose within the target range and to support a healthy pregnancy and birth.

\begin{itemize}
    \item \textbf{Healthy eating:} a dietitian helps plan meals that control blood glucose while providing good nutrition; the focus is on portion control, low glycaemic index carbohydrates, and regular meals
    \item \textbf{Physical activity:} regular, moderate activity such as walking helps lower blood glucose
    \item \textbf{Blood glucose monitoring:} checking levels at home and keeping a record for review
    \item \textbf{Insulin therapy:} needed by around 20--30\% of women when diet and activity alone do not control glucose
    \item \textbf{Oral medicines:} metformin is also used in some cases, either alone or with insulin
    \item \textbf{Monitoring the baby:} regular scans to check growth, and adjustments to birth planning if the baby is large
    \item \textbf{Birth planning:} most women with well-controlled GDM can have a vaginal birth at term; induction or caesarean may be recommended in some cases
    \item \textbf{Baby care after birth:} the newborn's blood glucose is checked after delivery
\end{itemize}

\section*{Complications}
\begin{itemize}
    \item Macrosomia: a large baby, which increases the risk of difficult birth, shoulder injury, and caesarean delivery
    \item Neonatal hypoglycaemia: low blood glucose in the baby after birth, which can occur if the baby produced excess insulin
    \item Pre-eclampsia and high blood pressure in the mother
    \item Increased risk of stillbirth, though this is rare with good management
    \item Respiratory distress syndrome and jaundice in the newborn
    \item Higher risk of obesity and type 2 diabetes in the child later in life
    \item Development of type 2 diabetes in the mother: around half of women with GDM develop it within 10 years
\end{itemize}

\section*{Prognosis and Outcomes}
With good management, the outlook for women with gestational diabetes and their babies is excellent, and most women have healthy pregnancies and healthy babies. Blood glucose control usually returns to normal soon after birth, but the diagnosis is an important warning sign: women with GDM face a high lifetime risk of type 2 diabetes. The postnatal period is therefore a crucial opportunity to adopt healthy eating, physical activity, and weight management, and to have regular diabetes screening. Babies of well-managed pregnancies have outcomes similar to those of unaffected pregnancies.

\section*{Prevention and Self-Care}
\begin{itemize}
    \item Maintain a healthy weight before pregnancy and gain weight within the recommended range during pregnancy
    \item Eat a balanced diet with plenty of vegetables, whole grains, and lean protein, and limit sugary drinks and refined carbohydrates
    \item Stay physically active before and during pregnancy
    \item Attend all antenatal appointments, including glucose screening at 24--28 weeks
    \item If diagnosed, follow the dietitian's and doctor's advice and monitor blood glucose as directed
    \item After the birth, have the postnatal glucose test and continue healthy habits
    \item Discuss your GDM history at future GP visits, as you need regular type 2 diabetes screening
\end{itemize}

\section*{Sources and Further Reading}
The following reputable sources provide further information for healthcare students and patients seeking reliable, current advice about gestational diabetes.

\begin{itemize}
    \item Diabetes Australia. \textit{Gestational diabetes fact sheet.} https://www.diabetesaustralia.com.au/
    \item Healthdirect Australia. \textit{Gestational diabetes.} https://www.healthdirect.gov.au/
    \item Jean Hailes for Women's Health. \textit{Gestational diabetes information.} https://www.jeanhailes.org.au/
    \item NHS. \textit{Gestational diabetes: overview and management.} https://www.nhs.uk/
    \item Mayo Clinic. \textit{Gestational diabetes: symptoms and causes.} https://www.mayoclinic.org/
    \item International Diabetes Federation. \textit{Gestational diabetes information.} https://www.idf.org/
\end{itemize}
